\documentclass[a4paper,12pt]{article}
\usepackage[T2A]{fontenc}
\usepackage[utf8]{inputenc}
\usepackage[russian]{babel}
\usepackage{amsmath,amssymb}
\usepackage{helvet}
\renewcommand{\familydefault}{\sfdefault}
\usepackage{icomma}
\usepackage{geometry}
\geometry{left=23mm,right=23mm,top=22mm,bottom=23mm}
\usepackage{microtype}
\usepackage{tikz}
\usetikzlibrary{arrows.meta,positioning,shapes.geometric,calc}
\usepackage{tkz-euclide}
\usepackage{tabularx,booktabs}
\usepackage{enumitem}
\usepackage{fancyhdr}
\usepackage{setspace}
\usepackage{xcolor}
\definecolor{accent}{HTML}{008080}
\definecolor{darkaccent}{HTML}{005C5C}
\definecolor{textgray}{HTML}{353B3D}
\definecolor{softgray}{HTML}{687276}
\setlength{\parindent}{0pt}
\setlength{\parskip}{6pt}
\onehalfspacing
\setlist[enumerate]{leftmargin=*,itemsep=10pt,topsep=8pt}
\pagestyle{fancy}
\fancyhf{}
\fancyhead[L]{\footnotesize\color{softgray}Медиана, биссектриса и высота}
\fancyhead[R]{\footnotesize\color{softgray}Кирилл Зиновьев}
\fancyfoot[C]{\footnotesize\color{softgray}\thepage}
\renewcommand{\headrulewidth}{0.4pt}
\renewcommand{\headrule}{\hbox to\headwidth{\color{accent}\leaders\hrule height \headrulewidth\hfill}}
\setlength{\headheight}{23pt}
\newcommand{\sectionline}{\par\vspace{-3mm}{\color{accent}\hrule height .5pt}\vspace{4mm}}
\newcommand{\key}[1]{\textcolor{darkaccent}{\textbf{#1}}}
\newcommand{\mathnote}[1]{{\small\color{softgray}#1}}
\begin{document}
\begin{titlepage}
\thispagestyle{empty}
\begin{center}
\vspace*{22mm}
{\color{darkaccent}\large МАТЕМАТИКА}\par
\vspace{15mm}
{\LARGE Чек-лист}\par
\vspace{13mm}
{\color{accent}\LARGE Медиана, биссектриса\\[2mm] и высота треугольника}\par
\vspace{13mm}
{\normalsize Геометрия. 7 класс}\par
\vfill
{\large Кирилл Зиновьев}\par
\vspace{12mm}
{28 октября 2026 г.}\par
\vspace{20mm}
{\small Лёвин Артём Александрович\\ эксклюзивно для Кирилла Зиновьева}\par
\vspace{14mm}
\begin{tikzpicture}[x=1cm,y=1cm]
\draw[accent,line width=.8pt] (0,0) .. controls (2.0,0.8) and (3.1,-0.7) .. (5.0,0)
  .. controls (6.9,0.7) and (8.0,-0.7) .. (10.0,0);
\end{tikzpicture}
\end{center}
\end{titlepage}
\setcounter{page}{1}

\section*{1. Три важных отрезка в треугольнике}
\sectionline
Из каждой вершины треугольника можно провести медиану, биссектрису и высоту. Для каждого отрезка есть свой \key{определяющий признак}.

\subsection*{1.1. Медиана делит сторону пополам}
\key{Медиана} --- отрезок, соединяющий вершину треугольника с серединой противоположной стороны.

\begin{center}
\begin{tikzpicture}[scale=.83]
\tkzDefPoint(1.6,3){A}
\tkzDefPoint(0,0){B}
\tkzDefPoint(5,0){C}
\tkzDefPoint(2.5,0){M}
\tkzDrawPolygon[thick](A,B,C)
\tkzDrawSegment[accent,very thick](A,M)
\tkzDrawPoints(A,B,C,M)
\tkzLabelPoints[below](B,M,C)
\tkzLabelPoint[above](A){$A$}
\tkzMarkSegments[mark=|,size=4pt](B,M M,C)
\end{tikzpicture}
\end{center}
Если $AM$ --- медиана, то $BM=MC$ и
\[
BM=MC=\frac{BC}{2},\qquad BC=2BM=2MC.
\]
\key{Пример.} $BC=18$ см. Тогда $BM=MC=18:2=9$ см. Если $BM=7{,}5$ см, то $BC=2\cdot7{,}5=15$ см.

\subsection*{1.2. Биссектриса делит угол пополам}
\key{Биссектриса угла треугольника} --- отрезок биссектрисы угла, соединяющий вершину с точкой противоположной стороны.

\begin{center}
\begin{tikzpicture}[scale=.82]
\tkzDefPoint(2,3){A}
\tkzDefPoint(0,0){B}
\tkzDefPoint(5,0){C}
\tkzDefPoint(2.30,0){D}
\tkzDrawPolygon[thick](A,B,C)
\tkzDrawSegment[accent,very thick](A,D)
\tkzDrawPoints(A,B,C,D)
\tkzLabelPoints[below](B,D,C)
\tkzLabelPoint[above](A){$A$}
\tkzMarkAngle[size=.70](B,A,D)
\tkzMarkAngle[size=.70](D,A,C)
\end{tikzpicture}
\end{center}
Если $AD$ --- биссектриса угла $BAC$, то
\[
\angle BAD=\angle DAC=\frac{\angle BAC}{2}.
\]
\key{Пример.} Если $\angle BAC=94^\circ$, каждая половина равна $47^\circ$. Если $\angle DAC=27^\circ$, весь угол $BAC$ равен $54^\circ$.

\clearpage
\section*{2. Высота и прямой угол}
\sectionline
\key{Высота треугольника} --- перпендикуляр, проведённый из вершины к прямой, содержащей противоположную сторону. Часть высоты между вершиной и этой прямой --- отрезок высоты.

\begin{center}
\begin{tikzpicture}[scale=.82]
\tkzDefPoint(1.6,3){A}
\tkzDefPoint(0,0){B}
\tkzDefPoint(5,0){C}
\tkzDefPoint(1.6,0){H}
\tkzDrawPolygon[thick](A,B,C)
\tkzDrawSegment[accent,very thick](A,H)
\tkzMarkRightAngle[size=.30](A,H,C)
\tkzDrawPoints(A,B,C,H)
\tkzLabelPoints[below](B,H,C)
\tkzLabelPoint[above](A){$A$}
\end{tikzpicture}
\end{center}
Если $AH$ --- высота и $H$ лежит на стороне $BC$, то
\[
AH\perp BC,\qquad \angle AHB=\angle AHC=90^\circ.
\]
\key{Пример.} Если $CH$ --- высота к стороне $AB$, причём $H$ находится между $A$ и $B$, то $\angle CHA=\angle CHB=90^\circ$.

\subsection*{Высота может опускаться на продолжение стороны}
В тупоугольном треугольнике основание высоты иногда лежит \key{за пределами стороны}. Для построения продолжите сторону до прямой.

\begin{center}
\begin{tikzpicture}[scale=.8]
\tkzDefPoint(-1.25,2.8){A}
\tkzDefPoint(0,0){B}
\tkzDefPoint(4.6,0){C}
\tkzDefPoint(-1.25,0){H}
\tkzDrawPolygon[thick](A,B,C)
\tkzDrawSegment[dashed,softgray](H,B)
\tkzDrawSegment[accent,very thick](A,H)
\tkzMarkRightAngle[size=.33](A,H,B)
\tkzDrawPoints(A,B,C,H)
\tkzLabelPoints[below](H,B,C)
\tkzLabelPoint[above](A){$A$}
\end{tikzpicture}
\end{center}
Здесь $AH\perp BC$, причём $H$ находится на \textbf{продолжении} $CB$ за точку $B$. Отрезок $AH$ всё равно является высотой.

\mathnote{Для точного построения используйте угольник: совместите одну его сторону с прямой, содержащей противоположную сторону, и проведите перпендикуляр через нужную вершину.}

\clearpage
\section*{3. Как читать и оформлять чертёж}
\sectionline
\begin{itemize}[leftmargin=*,itemsep=5pt]
\item \key{Равные отрезки} помечают одинаковым количеством штрихов. Одинаковые штрихи означают равенство длин.
\item \key{Равные углы} помечают одинаковыми дугами. Одна дуга и две дуги обозначают разные группы углов; их равенство из такой маркировки не следует.
\item \key{Прямой угол} отмечают маленьким квадратиком и при необходимости подписывают $90^\circ$.
\item Подписывайте \key{вершины}, точки на сторонах, основания высот и точки пересечения отрезков.
\end{itemize}

\subsection*{Три линии --- три точки пересечения}
В любом треугольнике:
\begin{enumerate}[label=\arabic*.,itemsep=3pt]
\item все \textbf{три медианы} пересекаются в одной точке;
\item все \textbf{три биссектрисы} пересекаются в одной точке;
\item прямые, на которых лежат \textbf{три высоты}, пересекаются в одной точке.
\end{enumerate}
Эти три точки в общем случае \textbf{различны}, хотя иногда совпадают. У тупоугольного треугольника точка пересечения прямых высот находится вне треугольника.

\subsection*{Если нужно вычислить угол}
Помните: сумма внутренних углов треугольника равна $180^\circ$:
\[
\angle A+\angle B+\angle C=180^\circ.
\]
Это правило помогает сначала найти целый угол, а затем разделить его биссектрисой пополам.

\subsection*{На что обращать внимание}
\begin{itemize}[leftmargin=*,itemsep=5pt]
\item Медиана гарантирует равенство \emph{частей стороны}, но не обязательно углов.
\item Биссектриса гарантирует равенство \emph{частей угла}, но не обязательно частей стороны.
\item Высота гарантирует угол $90^\circ$, но не обязательно делит сторону пополам.
\end{itemize}

\clearpage
\thispagestyle{fancy}
\begin{center}
{\Large\color{darkaccent}Алгоритм: как провести высоту}\par
\vspace{6mm}
\mathnote{От выбранной вершины к прямой, содержащей противоположную сторону}
\end{center}
\vspace{10mm}
\begin{center}
\begin{tikzpicture}[
  start/.style={ellipse,draw=accent,line width=.7pt,align=center,text width=6.6cm,minimum height=1.15cm},
  flowstep/.style={rectangle,rounded corners=2pt,draw=accent!70!black,line width=.6pt,align=center,text width=10.0cm,minimum height=1.35cm,inner sep=6pt},
  arr/.style={-{Stealth[length=3mm]},line width=.75pt,color=darkaccent},
  node distance=1.05cm]
\node[start] (s) {Начало};
\node[flowstep,below=of s] (one) {Выберите вершину треугольника};
\node[flowstep,below=of one] (two) {Найдите противоположную сторону и проведите прямую, содержащую её};
\node[flowstep,below=of two] (three) {Приложите одну сторону угольника к этой прямой};
\node[flowstep,below=of three] (four) {Передвиньте угольник, пока его перпендикулярная сторона не пройдёт через выбранную вершину};
\node[flowstep,below=of four] (five) {Проведите отрезок до прямой; при необходимости продолжите сторону треугольника};
\node[flowstep,below=of five] (six) {Отметьте основание высоты и прямой угол $90^\circ$};
\node[start,below=of six] (end) {Высота построена};
\foreach \from/\to in {s/one,one/two,two/three,three/four,four/five,five/six,six/end}{\draw[arr] (\from.south)--(\to.north);}
\end{tikzpicture}
\end{center}

\clearpage
\section*{4. Тренировочный блок --- Путь А}
\sectionline
\textit{10 заданий: от распознавания признака к задачам с несколькими шагами. Там, где требуется построение, используйте линейку и угольник; все точки подпишите.}
\begin{enumerate}[label=\textcolor{accent}{\bfseries\arabic*.}]
\item Назовите отрезок, проведённый из вершины треугольника, если он: а) делит противоположную сторону пополам; б) делит угол при вершине пополам; в) перпендикулярен прямой, содержащей противоположную сторону.
\item В треугольнике $ABC$ отрезок $AM$ --- медиана. Найдите $BM$ и $MC$, если $BC=26$ см.
\item В треугольнике $ABC$ отрезок $AM$ --- медиана и $BM=8{,}5$ см. Найдите $MC$ и $BC$.
\item В треугольнике $ABC$ отрезок $BD$ --- биссектриса угла $ABC$. Найдите $\angle ABD$ и $\angle DBC$, если $\angle ABC=106^\circ$.
\item В треугольнике $ABC$ отрезок $AD$ --- биссектриса угла $BAC$. Найдите $\angle BAD$ и $\angle BAC$, если $\angle DAC=31^\circ$.
\item В треугольнике $ABC$ проведена высота $CH$ к стороне $AB$; точка $H$ расположена между $A$ и $B$. Найдите $\angle CHA$ и $\angle CHB$.
\item Постройте три отдельных остроугольных треугольника. В первом проведите все медианы, во втором --- все биссектрисы, в третьем --- все высоты. Отметьте признаки каждого вида отрезков, подпишите точки их пересечения.
\item В треугольнике $ABC$ углы $B$ и $C$ равны $62^\circ$ и $48^\circ$. Отрезок $AD$ --- биссектриса угла $A$. Найдите $\angle BAD$ и $\angle DAC$.
\item В треугольнике $ABC$ проведены медиана $AM$ и биссектриса $BD$. Известно: $BC=34$ см, $\angle ABC=82^\circ$. Найдите $BM$, $MC$, $\angle ABD$ и $\angle DBC$.
\item В треугольнике $ABC$ известно: $\angle BAC=110^\circ$ и $\angle ABC=45^\circ$. Высота $BH$ опущена на прямую $AC$. Определите, на каком продолжении стороны $AC$ лежит точка $H$, и найдите $\angle ABH$ и $\angle BCA$.
\end{enumerate}

\clearpage
\section*{5. Ответы к тренировочному блоку}
\sectionline
\renewcommand{\arraystretch}{1.42}
\begin{tabularx}{\linewidth}{@{}p{11mm}X@{}}
\toprule
\textbf{№}&\textbf{Ответ или критерий проверки}\\
\midrule
1 & а) медиана; б) биссектриса; в) высота.\\
2 & $BM=MC=13$ см.\\
3 & $MC=8{,}5$ см; $BC=17$ см.\\
4 & $\angle ABD=\angle DBC=53^\circ$.\\
5 & $\angle BAD=31^\circ$; $\angle BAC=62^\circ$.\\
6 & $\angle CHA=\angle CHB=90^\circ$.\\
7 & На каждом из трёх чертежей проведены ровно три требуемых отрезка; у каждой тройки своя общая точка пересечения. У медиан отмечены равные части сторон, у биссектрис --- равные углы, у высот --- прямые углы.\\
8 & $\angle BAD=\angle DAC=35^\circ$.\\
9 & $BM=MC=17$ см; $\angle ABD=\angle DBC=41^\circ$.\\
10 & $H$ лежит на продолжении $CA$ за точку $A$; $\angle ABH=20^\circ$, $\angle BCA=25^\circ$.\\
\bottomrule
\end{tabularx}
\vspace{12mm}
\subsection*{Итоговая памятка}
\begin{tabularx}{\linewidth}{@{}p{34mm}X@{}}
\toprule
\textbf{В условии} & \textbf{Что сразу можно записать}\\
\midrule
Медиана $AM$ & $BM=MC$.\\
Биссектриса $AD$ & $\angle BAD=\angle DAC$.\\
Высота $CH$ & $CH\perp AB$ (к прямой $AB$); угол при $H$ равен $90^\circ$.\\
Любой треугольник & Сумма внутренних углов равна $180^\circ$.\\
\bottomrule
\end{tabularx}
\end{document}
